% SPDX-License-Identifier: Apache-2.0
% covers: DcSnoop
\datasheet{DcSnoop}{The data cache's snoop: the lines a write into the DDR3 covers, named to each hart that did not write them before the write is answered.}

\dsfacts{\code{cpu/vreteno/src/dcsnoop.rs}}%
{\code{vreteno32::dcsnoop::DcSnoop<P1>}}%
{\code{//docs:vreteno}}

\subsection*{Function}

The snoop keeps a hart's data cache of the DDR3 true when another host
writes the DDR3 (issue 1275). It stands on the router's port to the
DDR3, in the write address, write data and write response channels.
Reads do not pass through it.

The write address and data pass on as they come. A write burst into
the DDR3 is noted at its address phase: its identifier, the port it
came from, the first cache line it falls in and how many lines it
covers. When the burst's response comes back, which is when the
controller has made the write good, the snoop names the burst's lines
one a cycle on \code{inv}, to the first hart, and on \code{inv1}, to
the second. It holds the response until the last line is named and a
cycle more, so that a host's answer means the hart's copy is gone.

A line is named by its index alone, bits 4 to 11 of the address. That
may take out a line of another page, which the source calls harmless:
it keeps the compare out of the snoop.

Each line is named to each hart that did not write it: on \code{inv}
unless the writer was port 0, on \code{inv1} unless it was port
\code{P1}. With \code{P1} zero there is one hart: its own writes are
not noted, and \code{inv1} names nothing. With a second hart (issue
1408) every write into the DDR3 is noted, the first hart's too.

Whom to tell is decided when the response comes, not at the address
phase. The second hart is told nothing while its \code{asleep1} input
is set, that is while it waits in \code{wfi}: it forgets its cache as
it wakes. A noted write that nobody is to be told of is not held for a
walk: its answer passes as any other does.

\subsection*{Parameters}

\begin{center}\footnotesize
\begin{tabular}{@{}lp{0.7\textwidth}@{}}
\toprule
Parameter & Meaning \\
\midrule
\code{P1} & The arbiter port of the second hart, or zero for none. The
default is zero. \\
\bottomrule
\end{tabular}
\end{center}

Everything else is fixed in the source: addresses and data are 32 bits,
identifiers five bits with the port in the top three, as the board's
arbiter puts it, and the DDR3 is the addresses whose bits 31 and 30 are
\code{01}, \code{0x4000\_0000} to \code{0x7fff\_ffff}.

\subsection*{The line on \code{inv} and \code{inv1}}

Each is a \code{U<9>}. Bit 8 is set in a cycle that names a line, and
bits 0 to 7 are the line's index. In every other cycle the word is
zero.

The tables below are generated for \code{DcSnoop<7>}, the board's.

\dstables{DcSnoop}

\subsection*{Behaviour}

\begin{itemize}
\item \textbf{The table.} Four bursts may be noted at once, in the
slots \code{v0} to \code{v3}, each with its \code{id}, \code{line},
\code{cnt} and \code{pt}. A write to be noted while all four are full
waits at its address phase. A write not to be noted passes even then.
\item \textbf{The count.} A burst covers
\code{((addr[3:2] + len) >> 2) + 1} lines of sixteen bytes, from the
line its first word is in.
\item \textbf{Write data.} The beats pass on when the DDR3's side has
room, whatever the table holds.
\item \textbf{A response.} It is matched by identifier alone against
the noted bursts. Each host's identifiers are its own and none is
reused while its burst is out, as the source states. A response for no
noted burst passes. A response for a noted burst with someone to tell
starts the walk and waits; one with nobody to tell passes and empties
its slot.
\item \textbf{The walk.} \code{walking} is set, \code{wline} holds the
next line and \code{wleft} how many are left; one line is named each
cycle. \code{wport} and \code{wtell1} hold whom the walk is for. The
last line empties the slot in \code{wslot} and sets \code{settle}, the
cycle after the walk in which the response still waits.
\item \textbf{Waiting.} While \code{walking} or \code{settle} is set, no
response passes, not even one for a burst that was not noted.
\end{itemize}

\subsection*{Verification}

\begin{itemize}
\item Two tests in \code{dcsnoop.rs} run \code{DcSnoop<7>}, in
\code{//cpu/vreteno:isa\_test}. Ports 0, 7 and 3 each write one word
of a line of their own.
\begin{itemize}
\item \code{each\_hart\_is\_told\_of\_the\_others\_writes}: hart 0 is
told the lines ports 7 and 3 wrote, hart 1 the lines ports 0 and 3
wrote, and all three answers pass.
\item \code{a\_waiting\_hart\_is\_told\_nothing} sets \code{asleep1}:
hart 1 is told nothing, and port 0's answer passes earlier than with
hart 1 awake.
\end{itemize}
\item \code{//cpu/vreteno:board\_test} runs the board.
\begin{itemize}
\item \code{the\_data\_cache\_sees\_another\_hosts\_write} has the JTAG
host write a word of the DDR3 that the core has cached, or is filling,
and the core must read the new word.
\item \code{two\_harts\_lose\_no\_count} has each hart hand the other a
word the other has cached.
\item \code{a\_hart\_woken\_as\_a\_write\_lands\_sees\_it} wakes hart 1
on either side of the cycle a write's answer comes.
\item Two time stores into the DDR3 with the snoop on the path:\\
\code{both\_harts\_storing\_share\_the\_path} and\\
\code{the\_ddr3\_path\_is\_timed\_by\_the\_core}.
\end{itemize}
\end{itemize}

\subsection*{Used by}

\code{Board}, as \code{dcsnoop}, a \code{DcSnoop<7>}: its write
channels are the router's port 3, the DDR3, and its \code{inv} and
\code{inv1} go to hart 0 and hart 1, the second hart on the arbiter's
port 7. Hart 1's \code{asleep} output is its \code{asleep1}.

\subsection*{Limits}

Four bursts at a time; a fifth to be noted waits. One walk at a time,
and every response waits while it runs. A line is named by index alone,
so a line of another page with the same index goes too. A burst's span
is read as words from \code{addr} and \code{len}; \code{size} and
\code{burst} are not looked at. The range of the DDR3, the identifier
width and the port's place in it are fixed in the source, not
parameters.
